\documentclass[10pt,a4paper]{article}
\usepackage[T1]{fontenc}\usepackage[utf8]{inputenc}
\usepackage{lmodern,amsmath,amssymb,booktabs,array,microtype}
\pdfmapfile{+lm.map}\pdfmapfile{+cm.map}\pdfmapfile{+cmextra.map}\pdfmapfile{+symbols.map}
\pdfinfoomitdate=1\pdftrailerid{}
\usepackage[a4paper,margin=23mm]{geometry}
\usepackage{xcolor,hyperref,fancyhdr}
\definecolor{ink}{HTML}{183A4A}
\hypersetup{colorlinks=true,linkcolor=ink,urlcolor=ink,pdftitle={CST-DUST: Obere Reduktionsfehlergrenze v0.1}}
\pagestyle{fancy}\fancyhf{}\fancyhead[L]{\small CST-DUST | Obere Reduktionsfehlergrenze}
\fancyhead[R]{\small v0.1 | 11.11.2026}\fancyfoot[C]{\thepage}
\setlength{\headheight}{14pt}\setlength{\parindent}{0pt}\setlength{\parskip}{5pt}
\setlength{\emergencystretch}{2em}
\newcommand{\dd}{\,\mathrm d}\newcommand{\ii}{\mathrm i}\newcommand{\R}{\mathbb R}
\newcommand{\one}{\mathbf1}\newcommand{\norm}[1]{\lVert#1\rVert}
\newcommand{\TT}{\mathcal T}\newcommand{\EE}{\mathbb E}\newcommand{\HS}{\mathrm{HS}}
\begin{document}
{\color{ink}\LARGE\bfseries Eine obere Reduktionsfehlergrenze}\par
{\large Alle vier Eingänge, vollständige Modenreste und bedingte Verklebung}\par
\textbf{Version 0.1, 11. November 2026.}\\
Projekt: \texttt{thomasmgfeller/x6-spectral-astrolabium}.\\
Eingefrorener Quellstand: \texttt{3b6a186749b313ea59f43292e43052451dc4941c}.

\section*{Ergebnis und Reichweite}
Für die ursprüngliche Familie von 33 Historien gibt es jetzt eine analytische
\emph{obere} Restschranke in der Testtopologie aus Schritt 1. Sie umfasst alle
vier Eingänge, ihre kohärenten Überlagerungen und eine größere, unten definierte
unendliche Zustandsklasse. Feld-, Dichte- und Uhrmoden werden vollständig
kontrolliert. Die rückwärts gerichteten Paarungen erhalten eigene Konstanten.

Für die Verklebung gilt eine zweite Schranke: Sie kontrolliert die gemeinsame
Summe über alle $33^2=1089$ minimal verklebten Historien auf dem vierdimensionalen
äußeren Prüfraum. Dazu bleiben die bedingten Schnittdaten aus dem erweiterten
Randvertrag erhalten. Zwei bereits integrierte Randkerne werden nicht als
gewöhnliche Hilbert-Operatoren multipliziert.

\begin{center}\begin{tabular}{p{.43\linewidth}p{.48\linewidth}}\toprule
Gegenstand & Zertifiziertes Beispiel bei Toleranz $\tau=0{,}01$\\\midrule
Einzelkern: vollständiger Rest auf $\mathcal C_{\sqrt{261}}$ &
$E_{\rm single}<0{,}009004$ in $V'$, in beiden Richtungen.\\
Verklebung: beide äußeren Zustände in $E_4$ &
$E_{\rm sewn}<0{,}000970$ als absoluter Amplitudenfehler, vor numerischer Koeffizientenauswertung.\\
Hinreichende Basisgrößen dieser Abschätzung &
$3{,}385\cdot10^{18}$ beim Einzelkern; $3{,}503\cdot10^{34}$ bedingte Schnittmoden.\\
Praktische Entscheidung bei $10^6$ Moden &
\textbf{HOLD}: Die Abschätzung zertifiziert hier keine kompakte Rechnung mit dieser Toleranz.\\\bottomrule
\end{tabular}\end{center}

\textbf{Was gerechnet wurde.} Exakte algebraische Momente aller 33 Historien,
1089 Gaußpaarungen samt Abelvergleich, zehn unabhängige Testgruppen und eine
explizite Liste von 516 Startmoden. Die sehr großen Projektionsmatrizen wurden
\emph{nicht} aufgebaut. Die genannten Größen sind hinreichende, konservative
Zertifikate; sie sind weder Minimalgrößen noch Untergrenzen des tatsächlichen
Reduktionsfehlers.

\textbf{Abbruchregel.} Erfolg nur bei analytischem Rest plus zertifiziertem
numerischem Fehler unter $\tau$. Bei Überschreiten des Modenbudgets lautet das
Ergebnis \texttt{HOLD\_RESOURCE\_CAP}. Kleine neu beobachtete Matrixelemente
ersetzen die Restschranke nicht.

\textbf{Physikalische Grenze.} Das Ergebnis betrifft das deklarierte endliche
Kontrollmodell. Es konstruiert weder eine $L^2$-Evolution noch die vollständige
vierdimensionale CST-Summe und beweist keine Singularitätsauflösung.
\newpage

\section{Unveränderter Kern, Zustandsklasse und Fehlerbegriff}
Es gelten $\beta=4/\sqrt6$, $\kappa=1$, Skalarmasse null und ursprüngliche
Geometriekopplung eins. Die lokalen Kausalsets besitzen Anfang $\{0,1\}$,
Inneres $\{2\}$ und Ende $\{3,4\}$. Die in \texttt{prior\_core.py} exakt
aufgezählten 33 Ordnungen und ihr Gewicht $1/33$ bleiben unverändert.
\begin{align}
 c(n)&=(1,-9,16,-8)_n\ (0\le n\le3),\qquad c(n)=0\ (n\ge4),\notag\\
 F_4(C)&=|C|-\sum_{i\prec j}c(|I(i,j)|),\notag\\
 S_C/\hbar&=F_4(C)-\tfrac12\sum_j r_j
 -\tfrac\beta4\sum_{i\prec j}c_{ij}
 \bigl[r_j(t_j-t_i)^2+(q_j-q_i)^2\bigr].                 \tag{1}
\end{align}
Die Maße sind $\dd\lambda(t)=\dd t/2$ auf $[-1,1]$,
$\dd\nu(r)=\tfrac12\delta_0+\tfrac12\one_{(0,1)}\dd r$ und
$\dd\mu_A=\tfrac12\prod_{j=1}^2\dd\lambda(t_j)\dd\nu(r_j)\dd q_j$.
Randfunktionen sind unter gleichzeitiger Vertauschung der beiden vollständigen
Punktdatensätze invariant. Das ursprüngliche $H=L^2_{\rm inv}(\mu_A)$ bleibt
als Norm der vier Eingänge bestehen. Schritt 1 legt zusätzlich fest:
\[
 w(q)=(1+|q|^2)^4,\qquad V=L^2_{\rm inv}(w\dd\mu_A),\qquad
 V'=L^2_{\rm inv}(w^{-1}\dd\mu_A).                         \tag{2}
\]
Die Anti-Dualpaarung ist $\langle g,h\rangle=\int\bar g h\dd\mu_A$.
Es gilt $\norm{h}_{V'}=\sup_{\norm{g}_V\le1}|\langle g,h\rangle|$.
Der Abel-definierte ursprüngliche Kern $K:V\to V'$ ist beschränkt; eine
Abbildung aller normierten $H$-Zustände nach $H$ ist bereits widerlegt.

Mit $q_s=(q_1+q_2)/\sqrt2$, $q_a=(q_1-q_2)/\sqrt2$ lauten die Eingänge
\[
 h_0=\pi^{-1/2}e^{-(q_s^2+q_a^2)/2},\quad h_1=\sqrt2q_s h_0,
 \quad d_0=2\one_{00},\quad
 d_+=\frac2{\sqrt3}(\one_{10}+\one_{01}+\one_{11}),\quad
 \psi_{s,n}=\sqrt2\,d_s h_n.                              \tag{3}
\]
Die Uhrfunktion ist eins. Reihenfolge: $(0,0),(0,1),(+,0),(+,1)$.
Sie sind $H$-orthonormal; ihre $V$-Grammatrix ist
$\operatorname{diag}(65,261,65,261)$. Definiere
\[
 \begin{aligned}
 E_4&=\operatorname{span}\{\psi_{0,0},\psi_{0,1},\psi_{+,0},\psi_{+,1}\},\qquad
 \mathcal E_4=\{f\in E_4:\norm f_H\le1\},\\
 \mathcal C_M&=\{f\in V:\norm f_H\le1,\ \norm f_V\le M\},\qquad M=\sqrt{261}.
 \end{aligned}
                                                                  \tag{4}
\]
Damit ist $\mathcal E_4\subset\mathcal C_M$. Einzelantworten werden im gesamten
$V'$ geprüft, nicht nur an vier Matrixzeilen. Bei der Verklebung wird
$\sup_{f,g\in\mathcal E_4}|\langle g,(K_{\rm joint}-K_{\rm approx})f\rangle|$
kontrolliert. Diese zweite Aussage ist \emph{keine} Behauptung über die
$V\to V'$-Norm des gemeinsamen Kerns auf allen Zuständen.

\textbf{Rückwärts} bezeichnet die Anti-Dualadjungierte: die zum vertauschten,
komplex konjugierten Kern gehörende Testpaarung. Gemeint ist weder eine
inverse Dynamik noch unitarisierte Zeitumkehr.
\newpage

\section{Vollständige, verschachtelte Moden und richtige Projektion}
Setze $a=w^{-1/2}=(1+|q|^2)^{-2}$ und $Uf=a^{-1}f$. Dann ist $U:V\to H$
eine Isometrie. Der transformierte Kern
\[
 \TT=aKa:H\to H,\qquad
 \norm\TT_{\HS}\le C_*:=\frac\pi6\sqrt{\frac{2\pi}{\beta}}
 <1{,}027062                                                     \tag{5}
\]
ist Hilbert--Schmidt: $|K|\le\sqrt{2\pi/\beta}$ und
$\int_{\R^2}a(q)^2\dd q=\pi/3$. Der Faktor $1/2$ je Randmaß ergibt $\pi/6$.
Die endgültige Historieninterferenz bleibt in $K$ enthalten; nur die
\emph{obere Abschätzung} verwendet die Dreiecksungleichung.

Verwendet werden normierte zweidimensionale Hermitefunktionen
$h_{n_s,n_a}$ mit $n_s+n_a\le N$, sowie an jedem Punkt
\[
 \rho_{-1}=\sqrt2\,\one_{\{0\}},\quad
 \rho_k=\sqrt{2(2k+1)}P_k(2r-1)\one_{(0,1)}\ (0\le k\le D),\qquad
 \chi_k=\sqrt{2k+1}P_k(t)\ (0\le k\le L).                 \tag{6}
\]
Das Dichteatom bleibt ein eigener Modus. Aus den Produkten, mit Faktor $\sqrt2$
für das Randmaß, entstehen durch Punktvertauschung invariant symmetrisierte
Moden $e_\alpha$. Die relative Partnerphase ist $(-1)^{n_a}$; bei identischen
Punktindizes ist nur gerades $n_a$ erlaubt. Die Testmoden und Antwortmoden sind
\[
 v_\alpha=a e_\alpha\in V,\qquad
 \widetilde v_\alpha=w v_\alpha=a^{-1}e_\alpha\in V'.       \tag{7}
\]
Sie sind bezüglich der Paarung biorthogonal. Mit $P$ als orthogonalem
Projektor in $H$ und $P_V=U^{-1}PU$ ergibt sich
\[
 \TT_{N D L}=P\TT P,\qquad
 K_{N D L}f=\sum_{\alpha,\gamma}
 \widetilde v_\alpha\langle v_\alpha,Kv_\gamma\rangle
                  (v_\gamma,f)_V.                       \tag{8}
\]
Eine Expansion der Antwort in denselben $v_\alpha$ wäre im Allgemeinen falsch.
Alle Paarungen zwischen im Testspan enthaltenen Zuständen bleiben durch (8)
exakt erhalten. Das ist eine neue, der Testtopologie angepasste Reduktion;
sie wird nicht mit dem früheren $H$-Projektor gleichgesetzt.

Für $F=(N+1)(N+2)/2$ und $J=(D+2)(L+1)$ gilt
\[
 d_{\rm inv}=\tfrac12\bigl[FJ^2+(\lfloor N/2\rfloor+1)J\bigr].       \tag{9}
\]
$U\psi_{s,n}$ ist ein Gaußpaket mal Polynom vom Grad $4+n$.
Daher sind alle vier Eingänge ab $N\ge5,D\ge0,L\ge0$ exakt enthalten.
Die früheren 15 Testzustände, einschließlich Feldgrad zwei, Dichtemasken und
quadratischer Uhrmode, sind ab $(6,0,2)$ enthalten: 516 invariante Moden.
\texttt{basis.py} liefert deren vollständige Labels sowie die exakte Einbettung
der vier Eingänge. Jede Erhöhung von $N,D,L$ ergänzt einen ganzen Sektor;
es werden keine unbekannten Restmoden stillschweigend weggelassen.
\newpage

\section{Satz über den vollständigen Einzelkernrest}
Auf den jeweiligen Variablen wirken die positiven selbstadjungierten
Standardoperatoren mit Hermite- beziehungsweise Legendre-Spektrum:
\[
 H_q=-\Delta_q+|q|^2,\quad
 A_r=-\partial_r[r(1-r)\partial_r],\quad
 A_t=-\partial_t[(1-t^2)\partial_t].                       \tag{10}
\]
$A_r$ wirkt auf dem Atom als null und auf dem offenen Stratum mit der
natürlichen Legendre-Realisierung. Die verschwindenden Endpunktfaktoren
entfernen die Randterme. Die ersten ausgelassenen Eigenwerte sind
\[
 \lambda_q=2N+4,\qquad\lambda_r=(D+1)(D+2),\qquad
 \lambda_t=(L+1)(L+2).                                   \tag{11}
\]
Für die Ausgangs- und Eingangsvariablen seien
$C_{q,\pm}\ge\norm{H_{q,\pm}\TT}_{\HS}$ und entsprechend
$C_{r,j,\pm},C_{t,j,\pm}$ zertifiziert. Setze
\[
 A_\pm^2=\frac{C_{q,\pm}^2}{\lambda_q^2}
       +\sum_{j=1}^2\frac{C_{r,j,\pm}^2}{\lambda_r^2}
       +\sum_{j=1}^2\frac{C_{t,j,\pm}^2}{\lambda_t^2}.
                                                                  \tag{12}
\]
Dann gelten für den \emph{gesamten} ausgelassenen Rest:
\begin{align}
 \norm{(1-P)\TT}_{\HS}&\le A_+, &
 \norm{\TT(1-P)}_{\HS}&\le A_-,\notag\\
 \norm{K-K_{N D L}}_{V\to V'}&\le
 B:=\min\{C_*,\sqrt{A_+^2+A_-^2}\},                       \tag{13}\\
 \sup_{f\in\mathcal C_M}\norm{(K-K_{N D L})f}_{V'}&\le MB.
                                                                  \tag{14}
\end{align}
Dieselbe Schranke (14) gilt für den adjungierten Reduktionsfehler.
Für $f\in\mathcal E_4$ im Projektionsraum verbessert sich die Vorwärtsschranke
zu $\norm f_V A_+$, die rückwärts gerichtete zu $\norm f_V A_-$.

\textbf{Beweis.} Der spektrale Satz gibt für jeden Sektor
$\norm{(1-P_q)u}^2\le\lambda_q^{-2}\norm{H_qu}^2$ und analog für die
vier kompakten Faktoren. Für die kommutierenden Tensorprojektoren gilt
$1-P\le\sum_j(1-P_j)$. Integration über den anderen Rand liefert (12).
Die Hilbert--Schmidt-orthogonale Zerlegung
\[
 \TT-P\TT P=(1-P)\TT+P\TT(1-P)
\]
ergibt (13); $P\TT P$ ist zudem die orthogonale Projektion im Kernraum,
also ist ihr Rest höchstens $\norm\TT_{\HS}$. Die Isometrien zwischen
$V,V'$ und $H$ geben (14). Alle verwendeten Differentialnormen sind endlich,
wie die exakten Momente auf der folgenden Seite zeigen. Damit handelt es
sich um eine Restschranke über unendlich viele Moden, nicht um eine endliche
Summe zusätzlich beobachteter Koeffizienten. \hfill$\square$

Der Beweis benutzt keine positive Historieninterferenz und keine
Hermitizität von $K$. Besonders die Dichtekonstanten unterscheiden sich
zwischen beiden Richtungen deutlich.
\newpage

\section{Exakte Momente und zertifizierte Konstanten}
Für jede Historie sei $H^{(0)}$ die rationale Skalarhessian aus (1),
$B=\{0,1,3,4\}$ und
\[
 Q_C=H^{(0)}_{BB}-H^{(0)}_{B2}(H^{(0)}_{22})^{-1}H^{(0)}_{2B},
 \quad \alpha_C=-\beta H^{(0)}_{22}\ge\beta.
\]
Nach der inneren Skalarintegration hat jeder Term den Betrag höchstens
$\sqrt{2\pi/\alpha_C}$. Für eine Feldrichtung schreibe den Phasengradienten
als $Dy+Ex$, mit den entsprechenden Blöcken aus $\beta Q_C$, und $G=EE^T$.
Das normierte Feldmoment ist exakt
\begin{align}
 J_q={}&\frac{583}{35}+\frac{11}{5}\operatorname{tr}D^2
 -\frac15(\operatorname{tr}D)^2+\frac{53}{20}\operatorname{tr}G\notag\\
 &+\frac18\left[(\operatorname{tr}D^2)^2+2\operatorname{tr}D^4
 +\operatorname{tr}G\operatorname{tr}D^2+2\operatorname{tr}(DGD)
 +(\operatorname{tr}G)^2+2\operatorname{tr}G^2\right].       \tag{15}
\end{align}
Unter der normierten Dichte $(3/\pi)(1+|y|^2)^{-4}$ gilt
\[
 \EE\frac{|y|^{2m}}{(1+|y|^2)^k}
 =\frac{3m!\,\Gamma(3+k-m)}{\Gamma(4+k)},\qquad m<3+k.    \tag{16}
\]
Insbesondere sind $\EE|y|^2=1/2$ und $\EE|y|^4=1$.
Zur Kontrolle erhält man vor der Integration, nach Division durch den
Kernfaktor, für $H_y$ den Realteil
$|y|^2+|Dy+Ex|^2+(8-16|y|^2)/(1+|y|^2)^2$ und den Imaginärteil
$-\operatorname{tr}D+8y\cdot(Dy+Ex)/(1+|y|^2)$. Quadrieren und (16)
ergeben (15); im Test wird dieselbe Integration unabhängig kartesisch ausgeführt.

Für $\kappa_j=\tfrac12+\tfrac\beta4\sum_{i\prec j}c_{ij}(t_j-t_i)^2$
und $s_j=\partial_{t_j}(S_{\rm dust}/\hbar)$,
$s_{jj}=\partial_{t_j}^2(S_{\rm dust}/\hbar)$ sind die kompakten Momente
\[
 J_{r,j}=\EE_t\left[\frac{\kappa_j^2}{6}+\frac{\kappa_j^4}{60}\right],\quad
 J_{t,j}=\EE_{t,r}\left[(1-t_j^2)^2s_j^4+
 \{2t_js_j-(1-t_j^2)s_{jj}\}^2\right].                    \tag{17}
\]
Alle Erwartungswerte sind exakte Polynommomente der deklarierten Maße,
einschließlich sämtlicher Dichtestrata. Jensen kontrolliert die innere
kompakte Integration. Für jeden der fünf Differentialoperatoren gilt
\[
 C_{\bullet,\pm}=\frac\pi6\,\frac1{33}
 \sum_C\sqrt{\frac{2\pi}{\alpha_C}}\sqrt{J_{\bullet,C,\pm}}.        \tag{18}
\]
Die folgenden Dezimalzahlen dienen der Lesbarkeit; die Dateien enthalten
192-Bit-Arb-Einschließungen, die für alle Entscheidungen verwendet werden.
\begin{center}\begin{tabular}{lrr}\toprule
Konstante & vorwärts & rückwärts\\\midrule
$C_q$ &76,470398552781&76,470398552781\\
$C_{r,1}=C_{r,2}$ &2,776486574555&0,196205289258\\
$C_{t,1}=C_{t,2}$ &6,685137856442&6,301860085436\\\bottomrule
\end{tabular}\end{center}
\newpage

\section{Alle vier Eingänge und ihr Reduktionsfehler}
Die Anfangspunkte haben in jeder Historie keine Vorgänger. Ihre Dichtephase
ist stets $e^{-\ii r/2}$. Mit
\[
 E=\int_0^1e^{-\ii r/2}\dd r=e^{-\ii/4}\operatorname{sinc}(1/4),\qquad
 \zeta=\frac{2E+E^2}{\sqrt3}
 =1{,}6033902781-0{,}5537886752\ii                         \tag{19}
\]
gilt \emph{für die volle Summe}, und sogar für jede einzelne Historie,
\[
 K\psi_{+,n}=\zeta K\psi_{0,n},\qquad
 K\frac{\zeta\psi_{0,n}-\psi_{+,n}}{\sqrt{1+|\zeta|^2}}=0,
 \quad n=0,1.                                           \tag{20}
\]
Die beiden Nullzustände sind $H$-normiert. Eine einheitliche relative
Fehlertoleranz mit Nenner $\norm{Kf}$ ist daher auf $\mathcal E_4$ nicht
sinnvoll. Der absolute Fehler aus (2) bleibt wohldefiniert. Die Beziehung
(20) darf nicht auf die rückwärts gerichteten Antworten übertragen werden.
Globale Feldparität trennt $n=0$ und $n=1$ orthogonal; gemischte Paarungen
verschwinden exakt, auch unter den verwendeten Projektionen.

Ein hinreichendes analytisches Zertifikat für $\tau=1/100$ ist
\[
 (N,D,L)=(151306,124,192),\qquad
 d_{\rm inv}=3384672308841108222.                         \tag{21}
\]
Die Zahlen in der folgenden Tabelle sind \emph{nach oben gerundet}:
\begin{center}\begin{tabular}{lrr}\toprule
Eingang & Vorwärtsrest in $V'$ & Rückwärtsrest in $V'$\\\midrule
$\psi_{0,0}$ &0,003512098&0,002802409\\
$\psi_{0,1}$ &0,007037690&0,005615587\\
$\psi_{+,0}$ &0,003512098&0,002802409\\
$\psi_{+,1}$ &0,007037690&0,005615587\\\midrule
$\sup_{f\in\mathcal E_4}$ &0,007037690&0,005615587\\
$\sup_{f\in\mathcal C_{\sqrt{261}}}$, vollständiger Projektionsfehler
 &\multicolumn{2}{c}{$<0{,}009003549$}\\\bottomrule
\end{tabular}\end{center}
Gleiche Zeilengrenzen bedeuten keine Gleichheit der Antworten. Die Grenzen
sind obere Hüllen; positive Eingänge können zusätzlich mit $|\zeta|$ aus den
Atomeingängen abgeschätzt werden. Für kleine Basen nutzt der Code außerdem
die triviale $L^1$-Eingangsschranke und wählt das kleinere zertifizierte Ergebnis.

Die systematische Erweiterungsfolge wird rein auf Konstantenebene geprüft:
\begin{center}\begin{tabular}{rrrrl}\toprule
$N$&$D$&$L$&$d_{\rm inv}$&Status bei $10^6$ Moden, $\tau=0{,}01$\\\midrule
6&0&2&516&weiter ergänzen\\
12&1&4&10290&weiter ergänzen\\
24&3&8&329355&weiter ergänzen\\
48&7&16&14339925&\texttt{HOLD\_RESOURCE\_CAP}\\\bottomrule
\end{tabular}\end{center}
Die Formel kann auch über das Ressourcenlimit hinaus ausgewertet werden;
eine entsprechende Basisrechnung wurde damit nicht ausgeführt. Bereits
enthaltene vier oder 15 Matrixelemente allein begrenzen keine vollständige
Antwortnorm.
\newpage

\section{Verklebung: bedingte Paarung und Schnittgedächtnis}
Es gelten $L=\{0,1,2,3,4\}$, $R=\{3,4,5,6,7\}$ und $S=\{3,4\}$.
Der gemeinsame Vertrag aus Schritt 2 rekonstruiert aus Schnittinzidenzen,
Intervallprofilen und den benötigten früheren Feldern exakt
\[
 S_C=S_L+S_R-S_S+\Delta S.                                 \tag{22}
\]
Die fünf im vollständig geschichteten Beispiel fehlenden Terme und alle
anderen Kreuzterme in den 1089 minimalen Verklebungen bleiben enthalten.
$\Delta S$ hängt bei diesen minimalen Verklebungen nur von Endpunkten
\emph{außerhalb} des Schnitts ab. Der gesamte nicht am Schnitt liegende
Datensatz, einschließlich $q_2,q_5$, sei $\xi$. Für $\epsilon>0$ setze
\[
 A_\epsilon(s;\xi)=e^{-\epsilon|q_S|^2/4}
                    e^{\ii(S_L-S_S)/\hbar},\qquad
 B_\epsilon(s;\xi)=e^{-\epsilon|q_S|^2/4}
                    e^{-\ii S_R/\hbar}.                       \tag{23}
\]
Dann ist die Schnittintegration genau
\[
 2e^{\ii\Delta S/\hbar}\langle B_\epsilon,A_\epsilon\rangle_{H_S^{\rm lab}},
 \qquad H_S^{\rm lab}=L^2(\dd\mu_S),                        \tag{24}
\]
zusammen mit $e^{-\epsilon(q_2^2+q_5^2)/2}$ außerhalb des Schnitts.
Der Faktor zwei korrigiert das doppelte Schnittquotientenmaß, $S_S$ wird
einmal subtrahiert und der gemeinsame Regulator trifft jeden inneren Punkt
einmal. Das ursprüngliche Maß wird nicht durch $w^{-1}\dd\mu_S$ ersetzt.

\textbf{Warum eine gelabelte Schnittbasis?} Bei festgehaltenem $\xi$ und
festgehaltenen Historien sind $A_\epsilon,B_\epsilon$ unter Vertauschung
nur der Schnittpunkte im Allgemeinen \emph{nicht} invariant. Eine vorzeitige
Invarianzprojektion würde Information verwerfen. Hier wird daher die volle
Produktbasis (6) ohne Symmetrieeinschränkung verwendet:
\[
 d_{\rm lab}=FJ^2,\qquad \Pi=\Pi_{N D L}^{\rm lab}.
\]
Der Projektor kommutiert mit Relabeln; nach der gemeinsamen Historien- und
Datenintegration ist die äußere Paarung invariant. $\Pi$ ersetzt die
Schnittpaarung, während sämtliche Daten $\xi$ und $e^{\ii\Delta S/\hbar}$
erhalten bleiben.

Beide bedingten Funktionen sind für $\epsilon>0$ quadratintegrabel, mit
$\norm{A_\epsilon}^2=\norm{B_\epsilon}^2=\pi/\epsilon$. Die exakte Identität
\[
 \langle B,A\rangle-\langle\Pi B,\Pi A\rangle
   =\langle(1-\Pi)B,(1-\Pi)A\rangle                       \tag{25}
\]
zeigt, warum sowohl vorwärts als auch rückwärts gerichtete \emph{vollständige}
Modenreste gebraucht werden. Dies ist eine Hilfsrechnung mit bedingten,
regulierten Integranden. Sie setzt nicht voraus, dass $Kf\in H$ ist, und
liefert am Ende einen absoluten Fehler der zugelassenen Testpaarung.
\newpage

\section{Analytische obere Grenze beider bedingter Modenreste}
Am Schnitt sei die skalare Phase einer Seite $q^TDq/2+\ell^Tq$,
$\ell=\beta H^{(0)}_{S,\mathrm{rest}}q_{\rm rest}$. Setze
$d^2=\norm D_F^2$, $b^2=\norm{\beta H^{(0)}_{S,\mathrm{rest}}}_F^2$.
Unter der normierten Gaußdichte mit Kovarianz $\epsilon^{-1}I_2$ gilt
$\norm{|q|^2}_{L^2}=\sqrt8/\epsilon$. Direkte Differentiation von
$e^{-\epsilon|q|^2/4+\ii(q^TDq/2+\ell^Tq)}$, Dreiecksungleichung und
$\sqrt\epsilon|\ell|\le(\epsilon+|\ell|^2)/2$ ergeben
\begin{align}
 \frac{\norm{H_q A_\epsilon}}{\norm{A_\epsilon}}
 &\le c+3b^2|q_{\rm rest}|^2,\notag\\
 c&=\frac{3\epsilon}{2}+\frac{\sqrt8}{\epsilon}|1-\epsilon^2/4|
 +\frac{2\sqrt8}{\epsilon}d^2+\sqrt8 d+|\operatorname{tr}D|. \tag{26}
\end{align}
Dasselbe gilt für $B_\epsilon$. Eine Herleitung der Konstante: Der Realteil
von $H_qA/A$ ist $\epsilon+(1-\epsilon^2/4)|q|^2+|Dq+\ell|^2$;
der Imaginärteil ist $\epsilon q\cdot(Dq+\ell)-\operatorname{tr}D$.
Benutze $|Dq+\ell|^2\le2d^2|q|^2+2|\ell|^2$ sowie
$\sqrt\epsilon|\ell|\le\epsilon/2+|\ell|^2/2$. Die leicht vergrößerte
Konstante (26) folgt unmittelbar.

Die äußere Feldmaßhülle $|h_n|/\norm{h_n}_{L^1}$ hat Radialmomente
$n+2$ und $(n+2)(n+4)$. Der eine lokale innere Feldwert hat Gaußvarianz
$1/\epsilon$. Damit
\[
 m_2=n+2+\epsilon^{-1},\quad
 m_4=(n+2)(n+4)+2(n+2)\epsilon^{-1}+3\epsilon^{-2},\quad
 U_C=c^2+6cb^2m_2+9b^4m_4.                               \tag{27}
\]
Die Dichte- und Uhrenergien $R_C,T_C$ werden wie (17) exakt integriert.
Links wird wegen $S_L-S_S$ der konstante Anteil $1/2$ in $\kappa_j$
entfernt. Rechts ist $R_C=41/480$. Die rechten Uhrenergien werden für den
äußeren Dichtesektor $s=0,+$ getrennt berechnet: $00$ beziehungsweise das
gleichgewichtete Maß der Masken $10,01,11$. Definiere
\[
 L_C(n)=\sqrt{U_C/\lambda_q^2+R_C/\lambda_r^2+T_C/\lambda_t^2},
 \qquad R_C(s,n)\text{ analog für die rechte Seite}.        \tag{28}
\]
Die über die äußeren und lokalen inneren Daten gemittelte quadrierte
Schnittrestnorm ist höchstens $(\pi/\epsilon)L_C(n)^2$ beziehungsweise
$(\pi/\epsilon)R_C(s,n)^2$.

Seien $\bar L_n,\bar R_{s,n}$ die Mittelwerte von (28) über je 33 Historien,
$D_0=1/(2\sqrt2)$, $D_+=\sqrt3/(2\sqrt2)$ sowie
$\ell_0=\norm{h_0}_{L^1}=2\sqrt\pi$, $\ell_1=4$. Dann gilt
\[
 E^{\rm cut}_{s n,0 n}\le
 \frac{4\pi^2}{\epsilon^2}\ell_n^2 D_0D_s\,\bar L_n\bar R_{s,n}.       \tag{29}
\]
Beweis: (25), Cauchy--Schwarz, Jensen über die verbleibenden Daten und
$|e^{\ii\Delta S/\hbar}|=1$. Die äußeren Seitenmaße faktorisieren nach dem
Betragsübergang. Die zwei inneren Feldintegrale geben $2\pi/\epsilon$,
die Schnittnormen $\pi/\epsilon$, die Maßkorrektur zwei. Das erklärt
sämtliche Faktoren in (29), ohne einen Kreuzterm aus der Amplitude zu löschen.
\newpage

\section{Gemeinsamer Abelrest ohne Differenzheuristik}
In (29) sind positive Eingangsspalten exakt $\zeta$ mal die Atomspalten;
die globale Parität löscht die gemischten $n$-Blöcke. Somit ist auf
$\mathcal E_4\times\mathcal E_4$
\[
 E_{\rm cut}\le\max_{n=0,1}\sqrt{1+|\zeta|^2}
 \sqrt{(E^{\rm cut}_{0n,0n})^2+(E^{\rm cut}_{+n,0n})^2}.    \tag{30}
\]
Der äußere Zustandsraum ist damit vollständig abgedeckt, einschließlich
kohärenter Überlagerungen und Rückwärtspaarungen.

Es verbleibt $\epsilon\downarrow0$. Für jede gemeinsame Historie sei
$H^{(0)}_C$ die vollständige rationale $8\times8$-Hessian und
\[
 Q_\epsilon=\operatorname{diag}(1_{\rm außen},\epsilon_{\rm innen})
              -\ii\beta H_C^{(0)},\qquad
 F_{00}(\epsilon)=\frac{16\pi^3}{\prod_{k=1}^8\sqrt{p_k(\epsilon)}},\quad
 F_{11}(\epsilon)=F_{00}(\epsilon)e_{\rm in}^TQ_\epsilon^{-1}e_{\rm out}.
                                                                  \tag{31}
\]
Die Gaußelimination erfolgt mit den äußeren Indizes $(0,1,6,7)$ zuerst,
dann $(2,3,4,5)$. $p_k$ sind ihre Schurpivots. $e_{\rm in}$ hat Einsen
bei $0,1$, $e_{\rm out}$ bei $6,7$. Der Faktor $16\pi^3$ enthält die
Normierung $h_0h_0=\pi^{-1}$ und das achtdimensionale Gaußintegral.
Die Momente von $h_1h_1$ ergeben gerade die Inversenmatrix in (31).

Für $\epsilon>0$ sind alle Pivots akkretiv. Bei $\epsilon=0$ werden für
alle 1089 Fälle Nichtverschwindung und Abstand vom negativen reellen
Verzweigungsschnitt mit Arb eingeschlossen. Die Wurzeln sind also die
stetig fortgesetzten Abelzweige. Obwohl 32 rein innere Skalarblöcke singulär
sind, sind alle mit den äußeren Gaußpaketen versehenen $Q_0$ invertierbar.
Die endlichen Gaußpaarungen existieren damit analytisch; es wird keine
unkontrollierte punktweise Grenzwertannahme verwendet.

Setze $\delta_n(\epsilon)=1089^{-1}\sum_C|F_{nn,C}(\epsilon)-F_{nn,C}(0)|$.
Die verbleibende Staubphase wird im Betrag durch eins beschränkt, alle
Dichtestrata sind im Maß enthalten. Aus $D_0^2=1/8$ und
$D_0D_+=\sqrt3/8$ folgt
\[
 E_{\rm Abel}\le\frac{\sqrt{1+|\zeta|^2}}4
                         \max\{\delta_0,\delta_1\}.       \tag{32}
\]
Dies ist eine eingeschlossene Differenz zum \emph{analytisch definierten}
Abelwert (31), keine Differenz zweier Quadraturregeln als Fehlerersatz.
Die Gravitationsphase hat Betrag eins und ist für beide Regulatorwerte gleich.

\begin{center}\begin{tabular}{rrrr}\toprule
$\epsilon$&$\delta_0$ kleiner als&$\delta_1$ kleiner als&$E_{\rm Abel}$ kleiner als\\\midrule
$1/7$&0,023788252&0,041998890&0,020675504\\
$1/49$&0,003915235&0,007374456&0,003630348\\
$1/343$&0,000572131&0,001089731&0,000536461\\\bottomrule
\end{tabular}\end{center}
Zusätzlich reproduziert die neue Gaußauswertung die eingefrorene komplexe
1089-Historien-Summe im vollständigen Atomsektor. Das prüft Phase, Maß und
Wurzelzweig unabhängig von den neuen Restgrenzen.
\newpage

\section{Verklebungszertifikat und ausführbare Abbruchregel}
Bei $\epsilon=1/343$ liefert die explizite, konservative Budgetzuweisung
\[
 (N,D,L)=(76474292424,1510,2288),\qquad
 d_{\rm lab}=35026448697292180628528179868001600.            \tag{33}
\]
Die eingeschlossenen Obergrenzen aus (30) und (32) erfüllen
\[
 E_{\rm cut}<0{,}000432938,\qquad
 E_{\rm Abel}<0{,}000536461,\qquad
 E_{\rm total}<0{,}000969398+E_{\rm num}.                  \tag{34}
\]
Für $\tau=0{,}01$ bleibt also ein numerisches Budget größer als
$0{,}009030602$. $E_{\rm num}$ muss aus einer separat zertifizierten
Auswertung der tatsächlich implementierten Koeffizienten und verbleibenden
Integrale stammen. Ein solches riesiges Matrixverfahren wurde nicht
berechnet. (34) ist die Schranke für die mathematisch exakt projizierte
Verklebung. Der Code gibt diese Unterscheidung ausdrücklich aus.

\textbf{Algorithmus.} Beginne mit $(6,0,2)$ und einem erklärten Modenlimit,
hier $10^6$. Berechne aus exakten Momenten und eingeschlossenen Konstanten:
\begin{enumerate}
 \item Einzelkern: $E_{\rm an}=\sqrt{261}\min(C_*,\sqrt{A_+^2+A_-^2})$.
 \item Verklebung: zuerst $E_{\rm Abel}(\epsilon)\le\tau/4$
 nach (31)--(32) sicherstellen; dann $E_{\rm an}=E_{\rm cut}+E_{\rm Abel}$.
 Die im Prüfpaket dokumentierte Suche verwendet $\epsilon=7^{-k}$,
 höchstens elf Kandidaten; fehlende Zertifizierung führt zum Abbruch.
 \item Wenn Modenzahl größer als Limit: \texttt{HOLD\_RESOURCE\_CAP}.
 Sonst bei $E_{\rm an}+E_{\rm num}<\tau$: \texttt{TOLERANCE\_PASS}.
 Bei Gleichheit oder unentschiedener Einschließung wird weiter ergänzt.
 \item Andernfalls: \texttt{EXTEND\_BASIS}. Die dokumentierte verschachtelte
 Folge verdoppelt Feldgrad, erweitert die positiven Dichtegrade und
 die Uhrgrade gemeinsam. Alternativ können die größten Beiträge in
 (12) beziehungsweise (28) zuerst verringert werden.
\end{enumerate}
Die Ressourcenentscheidung hat Vorrang vor einem rein analytisch möglichen
Erfolg bei enormen Basisgrößen. Bei der Verklebung müssen Regulator und
Basis gekoppelt verfeinert werden: $\epsilon\to0$ bei festem Schnittprojektor
ist durch diese Abschätzung nicht kontrolliert.

\textbf{Bedeutung von HOLD.} Die vorliegende Hülle ist zu grob, um auf der
dokumentierten kleinen Erweiterungsfolge einen Fehler von $0{,}01$ zu
zertifizieren. Das beweist nicht, dass der tatsächliche Fehler groß ist,
oder dass keine effizientere Basis beziehungsweise schärfere Schranke
existiert. Die stärkste neue Aussage ist die explizite Kontrolle aller
bisher ausgelassenen Moden in einer festgelegten Topologie.

\textbf{Außerhalb des Vertrags.} Nicht erfasst sind zusätzliche globale
Historien, Schnittumgehungen, geänderte Maßfaktoren, andere Kopplungen,
unbegrenzte Dichte- oder Uhrintervalle und ein unendlicher Volumengrenzwert.
Die Verklebungsaussage ist auf $\mathcal E_4$ an beiden äußeren Rändern
beschränkt. Ein allgemeiner gemeinsamer Operatornormsatz bleibt offen.
\newpage

\section{Reproduktion, Quellen und unabhängige Prüfungen}
Das ZIP enthält Herleitung, PDF, ausführbaren Zertifikatsgenerator,
Modenkonstruktor, Budgetprüfer, sämtliche Resultate und acht bytegenau
übernommene Quelldateien. \texttt{provenance.json} fixiert deren Git-Blobs
auf dem eingangs genannten Main-Stand. \texttt{SHA256SUMS.json} schützt
alle mitgelieferten Dateien gegen versehentliche Änderungen.

Benötigt werden Python 3.12 und die im Paket fixierten Versionen von
SymPy, python-flint und NumPy. Die Zertifikate verwenden Arb mit 192 Bit;
NumPy dient nur unabhängigen Regressionskontrollen, nicht den beweisenden
Fehlerschranken. Aus einem entpackten Verzeichnis:
\begin{verbatim}
python -m pip install -r requirements.txt
python -B verify.py
python -B check_budget.py --field 24 --density 3 --clock 8
python -B check_budget.py --field 24 --density 3 --clock 8 \
    --composition
\end{verbatim}
\texttt{verify.py} prüft Quelldateien und Manifest, führt alle zehn Testgruppen
aus und erzeugt sämtliche sechs Resultatdateien in einer temporären Kopie
neu; anschließend vergleicht es sie bytegenau. Kein großes Modenarray und
keine große Galerkinmatrix werden dabei aufgebaut. Für ein selbst gewähltes
Numerikbudget gibt es \texttt{--numerical-error}; null bezeichnet ausschließlich
die mathematisch exakte Projektion.

\begin{center}\begin{tabular}{p{.09\linewidth}p{.81\linewidth}}\toprule
Test & Unabhängige Kontrolle\\\midrule
1 & Alle 66 Feldmomente durch kartesische Monomintegration statt Spurformel.\\
2 & Alle 132 Uhrmomente durch gradgenaue Tensor-Gaußregeln; reiner Regressionstest.\\
3 & Dichteableitungsnorm durch direkte Integration, Atom separat.\\
4 & Exakte Vierzustandseinbettung, vorherige Feld-/Uhrgrade, 516 Labels und Symmetrie.\\
5 & Fehlende Anfangsvorgänger in allen 33 und 1089 Historien; Dichtefaktor.\\
6 & Monotone Resthüllen, zertifizierte Gesamtbudgets und Ressourcenstatus.\\
7 & Gaußphase und ungerade Momente gegen unabhängige Eigenwert-/Inversenauswertung.\\
8 & Beide bedingten Restfaktoren und gelabelte statt vorzeitig symmetrisierter Schnittbasis.\\
9 & Abbruchentscheidungen einschließlich zusätzlichem numerischem Fehler.\\
10 & Gewichtete Grammnormen 65 und 261 durch unabhängige Radialintegration.\\\bottomrule
\end{tabular}\end{center}

\textbf{Primärquellen der Ausgangsdefinitionen.}
Die Dateien \texttt{inputs/prior\_domain\_DERIVATION.tex} und
\texttt{inputs/prior\_sewing.py} fixieren den Kernbereich und die exakte
Rekonstruktion der gemeinsamen Wirkung. Die früheren vier und 15 Testzustände
sind in den eingefrorenen Rand- und Reduktionsherleitungen dokumentiert.
Die Mathematik der orthogonalen Basisfunktionen folgt den Standarddefinitionen
im NIST Digital Library of Mathematical Functions,
\href{https://dlmf.nist.gov/18.3}{§18.3 (Orthogonalität)} und
\href{https://dlmf.nist.gov/18.8}{§18.8 (Differentialgleichungen)}.
Die konkreten Fehlerformeln, Konstanten und Prüfungen dieses Pakets werden
hier hergeleitet; sie werden nicht als Resultat dieser Referenz ausgegeben.

Lizenz des Prüfpakets: Apache-2.0, entsprechend dem Repository.
\end{document}
